\documentclass[a4paper]{arrowhead}

\usepackage[yyyymmdd]{datetime}
\usepackage{etoolbox}
\usepackage[T1]{fontenc}
\usepackage[utf8]{inputenc}
\usepackage{multirow}
\usepackage{longtable}
\usepackage{amssymb}
\usepackage{float}
\usepackage{tikz}
\usetikzlibrary{positioning,arrows.meta,fit,calc,shapes.misc,shapes.geometric}

\renewcommand{\dateseparator}{-}
\setlength{\emergencystretch}{3em}
\lstset{basicstyle=\ttfamily\footnotesize,breaklines=true,columns=fullflexible,frame=single,rulecolor=\color{gray!50},backgroundcolor=\color{gray!6}}
\tikzset{
  ahsys/.style={draw=ArrowheadBlue,very thick,fill=ArrowheadBlue!6,rounded corners=2pt,minimum width=3.2cm,minimum height=1.1cm,align=center,font=\small\bfseries},
  ahext/.style={draw=ArrowheadGrey,thick,dashed,fill=gray!5,rounded corners=2pt,minimum width=3.2cm,minimum height=1.0cm,align=center,font=\small},
  ahsrv/.style={draw=ArrowheadPurple,thick,fill=ArrowheadPurple!8,rounded corners=6pt,minimum height=0.55cm,inner xsep=4pt,align=center,font=\scriptsize\ttfamily},
  ahcsm/.style={draw=ArrowheadTeal,thick,fill=ArrowheadTeal!8,rounded corners=6pt,minimum height=0.55cm,inner xsep=4pt,align=center,font=\scriptsize\ttfamily},
  ahdb/.style={draw=ArrowheadOrange,thick,fill=ArrowheadOrange!8,cylinder,shape border rotate=90,aspect=0.25,minimum height=0.9cm,align=center,font=\scriptsize},
  ahflow/.style={-{Stealth[length=2.2mm]},thick,ArrowheadGrey},
  ahlabel/.style={font=\scriptsize,fill=white,inner sep=1pt,align=center},
}

%% Special references
\newcommand{\mref}[1]{{\textcolor{ArrowheadPurple}{\hyperref[sec:model:#1]{#1}}}}
\newcommand{\oref}[1]{{\textcolor{ArrowheadBlue}{\hyperref[sec:operations:#1]{#1}}}}
\newcommand{\pdef}[1]{{\textcolor{ArrowheadGrey}{#1\label{sec:model:primitives:#1}\label{sec:model:primitives:#1s}\label{sec:model:primitives:#1es}}}}
\newcommand{\pref}[1]{{\textcolor{ArrowheadGrey}{\hyperref[sec:model:primitives:#1]{#1}}}}

\newrobustcmd\osubsection[3]{
  \addtocounter{subsection}{1}
  \addcontentsline{toc}{subsection}{\protect\numberline{\thesubsection}operation \textcolor{ArrowheadBlue}{#1}}
  \subsection*{
    \thesubsection\quad
    operation
    \textcolor{ArrowheadBlue}{#1}
    (\notblank{#2}{\mref{#2}}{})
    \notblank{#3}{: \mref{#3}}{}
  }
  \label{sec:operations:#1}
}
\newrobustcmd\msubsection[2]{
  \addtocounter{subsection}{1}
  \addcontentsline{toc}{subsection}{\protect\numberline{\thesubsection}#1 \textcolor{ArrowheadPurple}{#2}}
  \subsection*{\thesubsection\quad#1 \textcolor{ArrowheadPurple}{#2}}
  \label{sec:model:#2} \label{sec:model:#2s} \label{sec:model:#2es}
}

\begin{document}

%% Arrowhead Document Properties
\ArrowheadTitle{sysMonitor}
\ArrowheadServiceID{sysMonitor}
\ArrowheadType{Service Description}
\ArrowheadTypeShort{SD}
\ArrowheadVersion{5.4.0}
\ArrowheadDate{2026-11-25}
\ArrowheadAuthor{Documentation maintenance}
\ArrowheadStatus{REVIEW DRAFT}
\ArrowheadContact{Repository maintainers}
\ArrowheadFooter{\href{www.arrowhead.eu}{www.arrowhead.eu}}
\ArrowheadSetup
%%

%% Front Page
\begin{center}
  \vspace*{1cm}
  \huge{\arrowtitle}

  \vspace*{0.2cm}
  \LARGE{\arrowtype}
  \vspace*{1cm}

  \Large{Service ID: \textit{``\arrowid''}}
  \vspace*{\fill}

  

  \vspace*{1cm}
  \vspace*{\fill}

  % Front Page Abstract
  \begin{abstract}
    This document describes the \texttt{sysMonitor} service of the Eclipse Arrowhead workflow support systems, which publishes observations of the WorkflowManager, here its own observation types about lost and suspected executors. It is provided by the WorkflowManager system, service version 1.0.0, documentation edition 5.4.0.
  \end{abstract}

  \vspace*{1cm}
\end{center}

\newpage
%%

%% Table of Contents
\tableofcontents
\newpage
%%

\section{Overview}
\label{sec:overview}
An executor can disappear while it holds work items. The WorkflowManager notices this only after some time, and a
bound step that is repeated may be done twice while the first executor still uses the equipment. The two observation
types of this document make both visible: \texttt{executor\allowbreak{}Liveness} and \texttt{workItemLost}. They are published next to
\texttt{echo}, \texttt{power} and \texttt{status}, which every Arrowhead system publishes and which the sysMonitor
documents of the common library describe.

The rest of this document is organized as follows.
In Section~\ref{sec:operations}, we describe the abstract message operations provided by the service.
In Section~\ref{sec:model}, we end the document by presenting the data types used by the mentioned operations.

\subsection{Significant Prior Art}
The \texttt{sysMonitor} service of Eclipse Arrowhead 5 and its AMQP interface (SysMonitor-AMQP-IDD-5.2.2).

\subsection{How This Service Is Meant to Be Used}
The types are listed in the service registration (metadata \texttt{observation\allowbreak{}Types} and the operations of the
interface), where SysMonitor finds them. A user of SysMonitor selects them in the subscription of the provider
WorkflowManager, on their own or together with the common types; \texttt{*} includes them. SysMonitor stores the
selected observations, offers them in its history and publishes them on MQTT with the observation type in the topic.

\subsection{Important Delimitations and Dependencies}
\label{sec:delimitations}
The observations report; they release no equipment and stop no executor. Whether the equipment is free
must be established at the equipment itself. Silence is not a sign of health: the absence of \texttt{echo} from the
WorkflowManager means that its executors are not being watched. Detection takes up to one check interval, the time-out
of an unanswered call (about 30\,s for a host that does not answer) and, for an unreachable executor, the grace period.

\paragraph{Authorization.} The broker permits a publisher only its own routing keys, identified by its certificate. SysMonitor receives
the offer by orchestration and needs a ConsumerAuthorization grant for the WorkflowManager's \texttt{sysMonitor}.

\newpage
\section{Service Interface}
\label{sec:operations}
This section describes the interface of the \texttt{sysMonitor} service. Each subsection names an abstract operation,
an input type and an output type, in that order. The input type is named inside parentheses, while the output type is
preceded by a colon. Input and output types are only denoted when accepted or returned, respectively, by the operation
in question. All abstract data types named in this section are defined in Section~\ref{sec:model}. The concrete
interface (AMQP/TLS/JSON) is specified in the AMQP IDD of the same name.

\begin{figure}[ht!]
  \centering
  \resizebox{0.92\textwidth}{!}{
\begin{tikzpicture}
\node[ahsys] (prov) at (4.6,0) {WorkflowManager\\[-1pt]{\normalfont\scriptsize $\ll$provider$\gg$}};
\node[ahsrv] (srv) at (1.6,0) {sysMonitor};
\draw[thick,ArrowheadPurple] (srv.east) -- (prov.west);
\node[ahext] (c0) at (-3.0,0.00) {SysMonitor\\[-1pt]{\scriptsize $\ll$consumer$\gg$}};
\draw[ahflow] (c0.east) -- (srv.west);
\end{tikzpicture}
  }
  \caption{SysML block view of the \texttt{sysMonitor} service: provider and typical consumers.}
  \label{fig:SD}
\end{figure}

The following interface operations are available.

{\small
\begin{longtable}{|p{4.4cm}|p{5.6cm}|p{5.8cm}|}

\hline
\rowcolor{gray!33} Operation & Input & Output \\ \hline
\endhead
\oref{executorLiveness} & -- & \mref{ExecutorLiveness} \\ \hline
\oref{workItemLost} & -- & \mref{WorkItemLost} \\ \hline
\end{longtable}
}

\osubsection{executorLiveness}{}{ExecutorLiveness}
State. It is published in every publication cycle (\texttt{sysmonitor.\allowbreak{}publish.\allowbreak{}interval}) and at once when an
executor becomes suspected, is cleared, refuses the question, or a work item is lost. A subscriber therefore learns of a
possible loss at the first failed check, before the grace period \texttt{lost\allowbreak{}After\allowbreak{}Seconds} has passed, and can secure the
equipment. Quality is \texttt{good} in the states \texttt{ok} and \texttt{disabled}, otherwise \texttt{uncertain};
\texttt{heldWorkItems} is as of the last check.
\osubsection{workItemLost}{}{WorkItemLost}
Event, one per lost work item, with quality \texttt{uncertain}. It is kept, and published again with the same
observation id, until the broker has confirmed it, so a subscriber may receive it more than once. At most 1000
unpublished events are kept, in memory: a restart of the WorkflowManager before the broker confirms an event loses the
event (the worklog of \texttt{workflow\allowbreak{}Management} still holds the loss).

\newpage
\section{Information Model}
\label{sec:model}
Here, all data objects that can be part of the \texttt{sysMonitor} service are listed in alphabetic order. Each
subsection, which describes one type of object, begins with the \textit{struct} keyword (a collection of named fields,
each with its own data type), the \textit{enum} keyword (a closed set of values) or the \textit{alias} keyword. As a
complement to the explicitly defined types in this section, there is also a list of implicit primitive types in
Section~\ref{sec:model:primitives}.

\msubsection{struct}{ExecutorLiveness}
Which executors are suspected of no longer running the work items they hold.

{\small
\begin{longtable}{|p{3.9cm}|p{4.4cm}|p{7.5cm}|}

\hline
\rowcolor{gray!33} Field & Type & Description \\ \hline
\endhead
\texttt{state} & \mref{LivenessState} & Summary of the last check. \\ \hline
\texttt{heldWorkItems} & \pref{Integer} & Work items dispatched to, claimed by or bound to executors. \\ \hline
\texttt{suspected\allowbreak{}Executors} & \pref{List}$<$\mref{SuspectedExecutor}$>$ & Executors that did not confirm a work item they hold. \\ \hline
\texttt{unchecked\allowbreak{}Executors} & \pref{List}$<$\pref{SystemName}$>$ & Executors that refuse \texttt{workflow\allowbreak{}Execution.\allowbreak{}get-\allowbreak{}status}; nothing is known about them. \\ \hline
\texttt{lostWorkItems} & \pref{Integer} & Work items lost since the WorkflowManager started. \\ \hline
\texttt{checkSeconds} & \pref{Duration} & \texttt{workflow.\allowbreak{}executor.\allowbreak{}liveness.\allowbreak{}check-\allowbreak{}seconds}. \\ \hline
\texttt{lost\allowbreak{}After\allowbreak{}Seconds} & \pref{Duration} & \texttt{workflow.\allowbreak{}executor.\allowbreak{}liveness.\allowbreak{}lost-\allowbreak{}after-\allowbreak{}seconds}. \\ \hline
\end{longtable}
}

\msubsection{enum}{LivenessState}
Summary of the liveness check.

{\small
\begin{longtable}{|p{3.9cm}|p{11.9cm}|}

\hline
\rowcolor{gray!33} Value & Meaning \\ \hline
\endhead
\texttt{ok} & No executor is suspected. \\ \hline
\texttt{suspect} & At least one executor is suspected. \\ \hline
\texttt{unchecked} & None is suspected, but at least one refuses the question. \\ \hline
\texttt{disabled} & The check is off (\texttt{check-seconds} is 0). \\ \hline
\end{longtable}
}

\msubsection{enum}{LossCause}
Why an executor is suspected.

{\small
\begin{longtable}{|p{3.9cm}|p{11.9cm}|}

\hline
\rowcolor{gray!33} Value & Meaning \\ \hline
\endhead
\texttt{notRunning} & It answers that it does not run the work item. \\ \hline
\texttt{unreachable} & It does not answer. \\ \hline
\texttt{unregistered} & It is no longer in the ServiceRegistry. \\ \hline
\end{longtable}
}

\msubsection{enum}{LossOutcome}
What follows for the step of a lost work item.

{\small
\begin{longtable}{|p{3.9cm}|p{11.9cm}|}

\hline
\rowcolor{gray!33} Value & Meaning \\ \hline
\endhead
\texttt{nextExecutor} & It was not bound: the next executor takes it, without an extension. \\ \hline
\texttt{rebind} & x-rebind: another executor runs the step. \\ \hline
\texttt{retry} & x-retry: the same executor runs it after the backoff. \\ \hline
\texttt{sessionFailed} & S-FailInvoke: the retry policy allows no further attempt. \\ \hline
\end{longtable}
}

\msubsection{struct}{SuspectedExecutor}
An executor that did not confirm work items it holds.

{\small
\begin{longtable}{|p{3.9cm}|p{4.4cm}|p{7.5cm}|}

\hline
\rowcolor{gray!33} Field & Type & Description \\ \hline
\endhead
\texttt{executor} & \pref{SystemName} & The executor. \\ \hline
\texttt{cause} & \mref{LossCause} & Why it is suspected. \\ \hline
\texttt{since} & \pref{DateTime} & When its first work item was suspected. \\ \hline
\texttt{workItemCount} & \pref{Integer} & Suspected work items of this executor. \\ \hline
\texttt{workItems} & \pref{List}$<$\mref{SuspectedWorkItem}$>$ & The work items; at most 100 are listed over all executors. \\ \hline
\end{longtable}
}

\msubsection{struct}{SuspectedWorkItem}
A suspected work item.

{\small
\begin{longtable}{|p{3.9cm}|p{4.4cm}|p{7.5cm}|}

\hline
\rowcolor{gray!33} Field & Type & Description \\ \hline
\endhead
\texttt{workItemId} & \pref{WorkItemId} & The work item. \\ \hline
\texttt{sessionId} & \pref{SessionId} & Its session. \\ \hline
\texttt{action} & \pref{Name} & Its action. \\ \hline
\texttt{step} & \pref{Name} & Its step. \\ \hline
\texttt{attempt} & \pref{Integer} & Attempt number of the step. \\ \hline
\texttt{bound} & \pref{Boolean} & Whether the executor reported BOUND, after which it uses the equipment. \\ \hline
\texttt{lastSeen} & \pref{DateTime} & When the executor last confirmed or reported the work item. \\ \hline
\end{longtable}
}

\msubsection{struct}{WorkItemLost}
A work item that the WorkflowManager gave up because its executor is gone.

{\small
\begin{longtable}{|p{3.9cm}|p{4.4cm}|p{7.5cm}|}

\hline
\rowcolor{gray!33} Field & Type & Description \\ \hline
\endhead
\texttt{workItemId} & \pref{WorkItemId} & The work item. \\ \hline
\texttt{sessionId} & \pref{SessionId} & Its session. \\ \hline
\texttt{recipe} & \pref{String} & Recipe as \texttt{name@version}. \\ \hline
\texttt{action} & \pref{Name} & Its action. \\ \hline
\texttt{step} & \pref{Name} & Its step. \\ \hline
\texttt{attempt} & \pref{Integer} & Attempt number of the step. \\ \hline
\texttt{executor} & \pref{SystemName} & The executor that held it. \\ \hline
\texttt{cause} & \mref{LossCause} & Why it was given up. \\ \hline
\texttt{detail} & \pref{String} & The cause in words. \\ \hline
\texttt{bound} & \pref{Boolean} & Whether it was bound. \\ \hline
\texttt{lastSeen} & \pref{DateTime} & When the executor last confirmed or reported it. \\ \hline
\texttt{suspected\allowbreak{}Since} & \pref{DateTime} & When the first check failed. \\ \hline
\texttt{lostAt} & \pref{DateTime} & When it was given up. \\ \hline
\texttt{detection\allowbreak{}Seconds} & \pref{Duration} & \texttt{lostAt} minus \texttt{lastSeen}: the longest time the loss can have gone unnoticed. \\ \hline
\texttt{equipment\allowbreak{}May\allowbreak{}Be\allowbreak{}In\allowbreak{}Use} & \pref{Boolean} & True for a bound work item: an executor that is only cut off may still be running the step. \\ \hline
\texttt{outcome} & \mref{LossOutcome} & What follows for the step. \\ \hline
\texttt{workRepeated} & \pref{Boolean} & Whether a bound step is run again, so that its work may be done twice. \\ \hline
\end{longtable}
}


\subsection{Primitives}
\label{sec:model:primitives}
Types and structures mentioned throughout this document that are assumed to be available to implementations of this
service. The concrete interpretation of each of these types must be provided by any IDD document claiming to implement
this service.

{\small
\begin{longtable}{|p{3.2cm}|p{12.6cm}|}

\hline
\rowcolor{gray!33} Type & Description \\ \hline
\endhead
\pdef{Boolean} & One out of \texttt{true} or \texttt{false}. \\ \hline
\pdef{DateTime} & Pinpoints a moment in time (UTC). \\ \hline
\pdef{Duration} & A length of time in seconds. \\ \hline
\pdef{Integer} & A whole number. \\ \hline
\pdef{List} & An ordered collection of items of the type given in angle brackets. \\ \hline
\pdef{Name} & An identifier of a recipe, action, step, parameter or resource: a letter followed by at most 62 letters, digits, \texttt{-} or \texttt{\_}. \\ \hline
\pdef{SessionId} & The identifier the WorkflowManager gives a session. \\ \hline
\pdef{String} & A chain of characters. \\ \hline
\pdef{SystemName} & The name of an Arrowhead system as it appears in its certificate and in the ServiceRegistry. \\ \hline
\pdef{WorkItemId} & The identifier of one attempt to run one step of one session. \\ \hline
\end{longtable}
}


\newpage
\bibliographystyle{IEEEtran}
\bibliography{bibliography}

\newpage

\section{Revision History}
\subsection{Amendments}

\noindent\begin{tabularx}{\textwidth}{| p{1cm} | p{2.2cm} | p{1.6cm} | X | p{3.6cm} |} \hline
\rowcolor{gray!33} No. & Date & Version & Subject of Amendments & Author \\ \hline
1 & 2026-11-25 & \arrowversion & Initial edition: service description of sysMonitor, written from the Java reference implementation (ah5-Workflow-java). & Documentation maintenance \\ \hline
\end{tabularx}

\subsection{Quality Assurance}

\noindent\begin{tabularx}{\textwidth}{| p{1cm} | p{2.2cm} | p{1.6cm} | X |} \hline
\rowcolor{gray!33} No. & Date & Version & Approved by \\ \hline
1 & & \arrowversion & Pending review \\ \hline
\end{tabularx}

\end{document}
